\subsection{Deep Reinforcement Learning}

Deep Reinforcement Learning (DRL) merges reinforcement learning's decision-making with the pattern recognition capabilities of neural networks, making it especially powerful for handling complex, high-dimensional data in environments like the Forex market. DRL empowers trading algorithms with adaptability, positioning them to potentially revolutionise trading strategies.

\subsubsection{Fundamentals and History}

The rise of DRL as a new field within machine learning gained momentum when it became clear that neural networks could greatly enhance the capabilities of reinforcement learning agents in dealing with and understanding complex environments. The groundbreaking work of DeepMind \cite{van_hasselt_deep_2016}, which demonstrated the success of the Double Deep Q-Network (DDQN) algorithm in mastering a range of Atari 2600 games using inputs, marked a significant breakthrough for DRL. This achievement not only validated the efficacy of DRL in mastering high-dimensional state spaces but also highlighted its potential applicability to financial markets. In this context, the challenge lies not only in interpreting large amounts of data but also in extracting strategic insights that can guide real-time trading decisions. This historical advance, along with ongoing research, has sparked interest in applying DRL to Forex trading, known to be a very volatile and dynamic environment.

\subsubsection{Deep Value-Based RL Algorithms: Deep Q-Networks}

Deep learning has transformed traditional value-based reinforcement learning methods, giving rise to what is known today as deep value-based methods. These advanced algorithms employ neural networks, known for their function approximation capabilities, to manage and interpret larger state spaces. Such an approach enables agents to learn and operate effectively in environments with complex, high-dimensional sensory input, a task that conventional value-based methods could not scale too efficiently. Neural Fitted Q-Learning was one of the early techniques that laid the foundation for deep value-based learning \cite{mnih_playing_2013}. It utilises neural networks to approximate the Q-function, which represents the expected return of a given state-action pair. By fitting a neural network to the Q-values obtained from completed episodes, it significantly enhanced the agent's ability to generalise across states in large and continuous spaces. However, the field of deep value-based methods truly came into prominence with the introduction of the Deep Q-Network (DQN) algorithm \cite{mnih_human-level_2015}. DQN integrated key innovations such as a separate target network for stabilising the learning targets and an experience replay mechanism to break the temporal correlations within training batches. This approach allowed DQN to effectively learn control policies directly from high-dimensional inputs, a significant jump forward in the application of RL. Following DQN, improvements have been made such as the Double Deep Q-Network (DDQN) to address limitations in the DQN architecture \cite{van_hasselt_deep_2016}. DDQN mitigates the overestimation bias of the action values by decoupling the action selection from the action evaluation. This is achieved by employing two networks: a primary network for action selection and a target network for evaluation, enhancing the stability and accuracy of the value estimation process. Although deep value-based methods are foundational and revolutionary, they will not be the focus of this research.

\subsubsection{Deep Policy-Based RL Algorithms: Deep Policy Gradients}

Deep policy-based RL algorithms excel in environments with continuous action spaces. Unlike value-based methods that assign a value to each potential action, policy-based approaches directly optimise the policy to maximise expected returns. This quality is invaluable in the domain of Forex trading, characterised by a vast and continuous action space that demands both precision and adaptability. Among these approaches, the Deep Deterministic Policy Gradient (DDPG) algorithm stands out as a model-free off-policy actor-critic algorithm specifically designed for continuous action spaces \cite{lillicrap_continuous_2015}. By combining the direct policy optimisation of policy-based methods with the evaluation power of value-based approaches, DDPG achieves a stable and robust learning process. The algorithm consists of an actor component for selecting actions and a critic component that evaluates the actions by estimating their Q-values. In Algorithm~\ref{Codes:DeepDeterministicPolicyGradient} we can observe how DDPG uses neural networks to approximate both the actor and critic functions. This enables handling complex input dimensions commonly encountered in Forex markets. The utilisation of a replay buffer for experience storage and the strategic sampling of mini-batches are key features that enhance learning stability by mitigating the correlation between sequential data points. Furthermore, DDPG incorporates target networks, which are gradually updated versions of the actor and critic, to solidify the learning targets and encourages gradual improvement.

\begin{algorithm}[htb!]
\caption{Deep Deterministic Policy Gradient \cite{lillicrap_continuous_2015}.}
\label{Codes:DeepDeterministicPolicyGradient}
\begin{algorithmic}
\REQUIRE an actor $\mu(s | \theta^{\mu})$, a critic $Q(s, a | \theta^{Q})$, a discount factor $\gamma$, a soft update rate $\tau \ll 1$ and a minibatch size $N$
\STATE Randomly initialise the critic parameters $\theta^{Q}$ and the actor parameters $\theta^{\mu}$
\STATE Initialise the target parameters $\theta^{Q'} \leftarrow \theta^{Q}$ and $\theta^{\mu'} \leftarrow \theta^{\mu}$, and the replay buffer $\mathcal{B}$
\FOR{episode $= 1, M$}
    \STATE Initialise a random process $\mathcal{N}$ for action exploration
    \STATE Receive the initial observation state $s_1$
    \FOR{$t = 1, T$}
        \STATE Select the action $a_t = \mu(s_t | \theta^{\mu}) + \mathcal{N}_t$ from the current policy and the exploration noise
        \STATE Execute the action $a_t$ and observe the reward $r_t$ and the new state $s_{t+1}$
        \STATE Store the transition $(s_t, a_t, r_t, s_{t+1})$ in $\mathcal{B}$
        \STATE Sample a random minibatch of $N$ transitions $(s_i, a_i, r_i, s_{i+1})$ from $\mathcal{B}$
        \STATE Set $y_i = r_i + \gamma Q'(s_{i+1}, \mu'(s_{i+1} | \theta^{\mu'}) | \theta^{Q'})$
        \STATE Update the critic by minimising the loss: $L = \frac{1}{N}\sum_i (y_i - Q(s_i, a_i | \theta^{Q}))^2$
        \STATE Update the actor policy using the sampled policy gradient:
        \STATE $\nabla_{\theta^{\mu}} J \approx \frac{1}{N}\sum_i \nabla_a Q(s, a | \theta^{Q})|_{s = s_i, a = \mu(s_i)} \nabla_{\theta^{\mu}} \mu(s | \theta^{\mu})|_{s_i}$
        \STATE Update the target networks:
        \STATE $\theta^{Q'} \leftarrow \tau \theta^{Q} + (1 - \tau) \theta^{Q'}$
        \STATE $\theta^{\mu'} \leftarrow \tau \theta^{\mu} + (1 - \tau) \theta^{\mu'}$
    \ENDFOR
\ENDFOR
\end{algorithmic}
\end{algorithm}

Although we cannot guarantee the convergence of the DDPG due to the partially observable nature of the Forex markets, its design for continuous policy adaptation presents a good approach to developing sophisticated trading strategies. The adaptability of DDPG to the evolving conditions of financial markets and its ability to identify complex patterns highlight its potential in algorithmic trading.